\documentclass[11pt,a4paper]{article}
\usepackage[T1]{fontenc}
\usepackage{lmodern}
\usepackage{amsmath,amssymb,amsthm,mathtools,bm,mathrsfs}
\usepackage[margin=25mm]{geometry}
\usepackage{microtype,booktabs,array}
\usepackage[hidelinks]{hyperref}
\hypersetup{pdftitle={The tree-level two-scalar spectrum in global AdS3},pdfsubject={Low-spin completion from local form factors and a crossed Casimir equation}}
\usepackage{enumitem}
\setlist{nosep,leftmargin=*}
\setlength{\parskip}{4pt}
\setlength{\parindent}{14pt}
\allowdisplaybreaks[2]
\numberwithin{equation}{section}
\newtheorem{theorem}{Theorem}[section]
\newtheorem{lemma}[theorem]{Lemma}
\newtheorem{proposition}[theorem]{Proposition}
\theoremstyle{remark}\newtheorem{remark}[theorem]{Remark}
\newcommand{\AdS}{\mathrm{AdS}}
\newcommand{\dd}{\mathrm d}
\newcommand{\tr}{\operatorname{tr}}
\newcommand{\Boxop}{\Box}
\newcommand{\kap}{\kappa}
\newcommand{\calC}{\mathscr C}
\newcommand{\calL}{\mathcal L}
\newcommand{\ket}[1]{|#1\rangle}
\newcommand{\bra}[1]{\langle#1|}
\newcommand{\norm}[1]{\lVert#1\rVert}
\newcommand{\poch}[2]{(#1)_{#2}}
\newcommand{\arxiv}[1]{\href{https://arxiv.org/abs/#1}{arXiv:#1}}
\title{The tree-level two-scalar spectrum in global $\AdS_3$\\[5pt]
\large Low-spin completion from local form factors and a crossed Casimir equation}
\author{}
\date{Research draft --- 12 September 2026}
\begin{document}
\maketitle
\vspace{-15pt}
\begin{abstract}
We derive closed leading gravitational binding coefficients for every radial level and every allowed spin of two identical, minimally coupled real scalars in global $\AdS_3$. The physical one-scalar mass is held fixed; no independent four-scalar contact interaction is included. The known spin-independent high-spin result is completed by explicit spin-zero and spin-two terms. The proof separates crossed exchange from annihilation. Normalized local primary form factors determine the annihilation contribution, which vanishes above spin two. Acting on crossed exchange with the graviton Casimir collapses it to a local interaction and produces a three-term difference equation for its primary coefficients. A positive, explicit $2\times2$ determinant proves that the known high-spin tail and this equation uniquely fix all lower spins. An independent circular canonical reduction, followed by representation-theoretic reconstruction, agrees with the result separately in both channels. The result concerns connected tree-level scalar-primary data; it does not compute a bare one-particle self-energy, higher-loop corrections, or a complete interacting Fock-space mixing matrix.
\end{abstract}

\section{Introduction and scope}

The leading gravitational binding energy of two particles is a dynamical quantity even when the one-particle spectrum has been fixed by a mass-renormalization condition. In global anti-de Sitter space, this distinction can be made directly in canonical perturbation theory: an interacting primary energy is a shifted lowest weight, while its symmetry descendants retain integer level spacings. The use of binding energies to extract anomalous dimensions is developed, for example, in Ref.~\cite{FS}.

For $\AdS_3$ graviton exchange, the high-spin double-trace anomalous dimensions are already known. Inverting crossed-channel stress-tensor exchange gives a result depending on the smaller chiral level, but not on the spin \cite{KSS,CH}. This does not license an unqualified extrapolation to the spins at or below the exchanged spin. Low-spin terms depend on the actual bulk interaction and its contact prescription. The more general need to complete large-spin exchange data at finite spin is illustrated in a different dimension in Ref.~\cite{ABP}; that reference is not a source for the $\AdS_3$ coefficients derived here.

The purpose of this paper is to complete the minimal Einstein--real-scalar answer at spins zero and two, for arbitrary radial excitation. There are two useful structural facts. First, the annihilation-channel stress tensor of a normalized two-scalar primary can be obtained from its value at the centre and the global highest-weight equations. Second, the crossed-channel Casimir equation has a unique low-spin completion once its source and its high-spin tail are specified. Neither statement requires extrapolating a finite table of energies.

We use ``spectrum'' below in the following precise perturbative sense. The computed quantity is the connected tree-level energy-shift operator in the scalar two-particle representation, equivalently the double-trace logarithmic data of the minimal four-scalar exchange amplitude. For an isolated scalar-primary branch, it is the leading physical binding energy. At an accidental degeneracy with additional physical sectors, a complete spectral interpretation requires the corresponding full projected Hamiltonian. We do not construct that larger matrix here. Boundary gravitons are not declared proper gauge.

\section{Theory, canonical normalization, and the result}
\label{sec:result}

Set the AdS radius to one, choose signature $(-,+,+)$, and write
\begin{equation}
 \kap^2=16\pi G,\qquad \mu=m^2=\Delta(\Delta-2),\qquad \Delta>1.
\end{equation}
The regulated action and its boundary prescription are
\begin{align}
 S_R={}&\frac1{\kap^2}\int_{M_R}\!\sqrt{-g}\,(R+2)
 -\frac12\int_{M_R}\!\sqrt{-g}\,\bigl((\nabla\phi)^2+\mu\phi^2\bigr)
 +\frac2{\kap^2}\int_{\Gamma_R}\!\sqrt{-\gamma}\,(K-1).
 \label{eq:action}
\end{align}
We keep a smooth centre, a fixed boundary-cylinder metric, Brown--Henneaux metric falloffs, and the source-free standard scalar condition $\phi=O(r^{-\Delta})$. The boundary time is the global coordinate in
\begin{equation}
 \dd s_0^2=-f\dd t^2+\frac{\dd r^2}{f}+r^2\dd\varphi^2,
 \qquad f=1+r^2,\qquad \varphi\sim\varphi+2\pi.
\end{equation}
The boundary action and canonical endpoint terms are retained in taking the covariant phase-space reduction; see Ref.~\cite{HW} for the general boundary framework. There is no independent $O(G)\phi^4$ or higher-derivative four-scalar coupling in \eqref{eq:action}.

All energies are gaps above the interacting vacuum. We define the physical mass by the lowest scalar gap, $E^{(1)}_{00}=\Delta$. AdS symmetry then gives
\begin{equation}
 E^{(1)}_{nj}=\Delta+2n+|j|.
 \label{eq:oneparticle}
\end{equation}
This is not a computation of a bare-to-renormalized mass relation. The one-body self-energy and its finite mass counterterm have been absorbed into the specified physical $\Delta$.

Introduce
\begin{equation}
 h=\Delta+n,\qquad C=h(h-1),\qquad
 U_n=4(\mu-2C)=-4\bigl[\Delta^2+2n(2\Delta+n-1)\bigr].
 \label{eq:U}
\end{equation}
A two-scalar primary has free energy $2\Delta+2n+|\ell|$. For identical real bosons, $\ell$ is even.

\begin{theorem}[Connected tree-level scalar-primary coefficients]
\label{thm:main}
For the minimal theory \eqref{eq:action}, with the reflecting, regular, AdS-invariant exchange prescription and physical mass convention above, write
\begin{equation}
 E^{(2)}_{n\ell}=2\Delta+2n+|\ell|+Gg_{n\ell}+O(G^2).
\end{equation}
The connected coefficients for $n\geq0$ and even $\ell$ are
\begin{align}
 g_{n0}&=U_n+
 \frac{2\left[15C^2-(12+10\mu)C-\mu^2+6\mu\right]}
 {(2h-3)(2h-1)(2h+1)},\label{eq:full0}\\
 g_{n,\pm2}&=U_n+
 \frac{(n+1)(n+2)(2\Delta+n-1)(2\Delta+n)}
 {(2h-1)(2h+1)(2h+3)},\label{eq:full2}\\
 g_{n\ell}&=U_n,\qquad |\ell|\geq4.\label{eq:fullhigh}
\end{align}
At $n=0$, $\Delta=3/2$, the uncontracted form of \eqref{eq:full0} has a removable $0/0$; its continuous value is $g_{00}=-9/2$. These formulas are fixed-level perturbative coefficients, not a statement of uniform validity as the level tends to infinity.
\end{theorem}

The proof uses the established high-spin stress-tensor result in Ref.~\cite{KSS}, restricted conservatively to spin strictly above two. The derivation below determines the low-spin completion from the minimal bulk action, rather than extending the inversion formula outside its stated domain.

It is useful to record a more diagnostic channel decomposition:
\begin{align}
 g^{\mathrm x}_{n0}&=U_n\frac{2h-2}{2h-1},&
 g^{\mathrm x}_{n\ell}&=U_n\quad (|\ell|\geq1),\label{eq:crossanswer}\\
 g^{\mathrm s}_{n0}&=-\frac{2(C+\mu)^2}{(2h-3)(2h-1)(2h+1)},&
 g^{\mathrm s}_{n,\pm2}&=\frac{(n+1)(n+2)(2\Delta+n-1)(2\Delta+n)}{(2h-1)(2h+1)(2h+3)},\label{eq:sanswer}\\
 g^{\mathrm s}_{n\ell}&=0\quad(|\ell|>2),&g_{n\ell}&=g^{\mathrm x}_{n\ell}+g^{\mathrm s}_{n\ell}.
\end{align}
Here $\mathrm x$ denotes the combined crossed contribution in the identical-boson matrix element. Equivalently, it is the exchange energy for two distinguishable equal-mass species, which provides auxiliary odd-spin data used in the proof. No odd-spin identical-boson primary is being added to the physical spectrum. Notice especially that the spin-zero crossed answer itself differs from $U_n$.

\section{Tree exchange and normalized local primary form factors}
\label{sec:formfactors}

\subsection{Integrating the metric response}

For a free scalar configuration let $q$ solve
\begin{equation}
 \mathcal E^{(1)}[q]=\frac12T,\qquad
 T_{\mu\nu}=\partial_\mu\phi\partial_\nu\phi
 -\frac12g_{\mu\nu}\bigl((\partial\phi)^2+\mu\phi^2\bigr),
 \qquad g=g_0+\kap^2q+\cdots.
\end{equation}
Quadratic homogeneity of the Einstein action yields
\begin{equation}
 S_{\mathrm{eff},4}=\frac{\kap^2}{4}\int\dd V\,q_{\mu\nu}T^{\mu\nu}.
 \label{eq:effaction}
\end{equation}
The factor $1/4$ includes the on-shell quadratic Einstein term; retaining only the matter coupling would double the answer. In canonical normal form the resonant quartic Hamiltonian is minus the time-averaged interaction Lagrangian. For a normalized annihilation form factor
$\tau_{\mu\nu}=\bra0:T_{\mu\nu}:\ket{P}$, this gives
\begin{equation}
 \gamma^{\mathrm s}=-\frac{\kap^2}{2}
 \int_\Sigma r\,\dd r\,\dd\varphi\,
 q[\tau]_{\mu\nu}\tau^{*\mu\nu}.
 \label{eq:annpairing}
\end{equation}
The factor of two combines the positive--positive and negative--negative source pairings. Normal ordering here specifies the connected tree interaction, not the unrenormalized one-body self-energy.

The pairing is unchanged by a bounded pure-diffeomorphism addition with vanishing radial flux: conservation turns the change into endpoints, whose time average vanishes. This fact is used for the particular response, not to quotient independent physical Brown--Henneaux excitations.

\subsection{Chiral primary coefficients and centre values}

Let a normalized scalar mode have chiral levels $(p,q)$, energy $\Delta+p+q$, and angular momentum $p-q$. The normalized coefficients of a chiral two-particle primary at level $K$ can be chosen real:
\begin{equation}
 v_p^{(K)}=(-1)^p
 \sqrt{\frac{\binom Kp(\Delta)_K^2}
 {(\Delta)_p(\Delta)_{K-p}(2\Delta+K-1)_K}},\qquad 0\leq p\leq K.
 \label{eq:CG}
\end{equation}
Their norm is one by the terminating Vandermonde identity. A distinguishable two-scalar primary at chiral levels $(K,L)$ has product coefficients $v_p^{(K)}v_q^{(L)}$. Identical-boson symmetrization retains $K+L$ even and multiplies a local vacuum-to-pair form factor by $\sqrt2$.

For $j\geq0$, the scalar modes have the centre expansion
\begin{equation}
 u_{q+j,q}(0,r,\varphi)=
 \frac{(-1)^q}{\sqrt{2\pi}\,j!}
 \sqrt{\frac{(q+j)!}{q!}(\Delta+q)_j}
 r^je^{ij\varphi}+O(r^{j+2}).
 \label{eq:centremode}
\end{equation}
In particular, a distinguishable spin-zero primary with $K=L=n$ obeys
\begin{equation}
 F_n=\bra0\phi_1\phi_2\ket{P_{n0}}
 =\frac{(-1)^n}{2\pi}f^{-h}e^{-2iht},\qquad
 \Box F_n=4C F_n.
 \label{eq:F}
\end{equation}
The coefficient at the centre is just $(-1)^n\sum_p(v_p^{(n)})^2/(2\pi)$. The identical-boson form factor is $F_n^{\mathrm{id}}=\sqrt2F_n$.

For later use, a spin-$s$ covariant primary tensor can be written
\begin{align}
 W^{(E,s)}&=e^{-iEt+is\varphi}r^sf^{-E/2}v^{\otimes s},&
 v&=\dd t-\dd\varphi+\frac{i\dd r}{rf},\qquad s=1,2.\label{eq:W}
\end{align}
Direct differentiation gives
\begin{equation}
 \nabla\cdot W^{(E,s)}=0,\qquad
 \Box W^{(E,s)}=[E(E-2)-s]W^{(E,s)}.
 \label{eq:Weigen}
\end{equation}
The rank-two tensor is traceless. Both global lowering Lie derivatives annihilate it. Its apparent polar-coordinate singularities cancel in Cartesian components. Moreover,
\begin{equation}
 \int_\Sigma r\,\dd r\,\dd\varphi\,
 W^{(E,s)}\cdot W^{(E,s)*}
 =\frac{\pi 2^s}{E+s-1}.
 \label{eq:Wnorm}
\end{equation}
This is a tensor contraction in the exchange pairing, not a Hilbert-space norm for arbitrary off-shell tensors.

For the distinguishable current $J_\mu=\phi_1\nabla_\mu\phi_2-\phi_2\nabla_\mu\phi_1$, only a spin-one primary form factor survives. At $(K,L)=(n+1,n)$ it has the form
\begin{equation}
 J_n=A_1 W^{(2h+1,1)},\qquad
 |A_1|^2=\frac{h(n+1)(2\Delta+n-1)}{2\pi^2(2h-1)}.
 \label{eq:Jnorm}
\end{equation}
For the identical scalar stress tensor at spin two,
\begin{equation}
 \tau_{n2}=A_2W^{(2h+2,2)},\qquad
 |A_2|^2=\frac{h(h+1)(n+1)(n+2)(2\Delta+n-1)(2\Delta+n)}
 {8\pi^2(2h-1)(2h+1)}.
 \label{eq:taunorm}
\end{equation}
The elementary coefficient ratios proving these expressions at arbitrary $n$ are given in Appendix~\ref{app:normalization}.

\subsection{Why the local form factors have finite spin support}

These are primary form factors, not arbitrary tensor waves of fixed angular momentum. They obey both first-order global lowering equations. Their analytic regular solution is determined by the tensor value at the centre: the two lowering vectors span the complexified spatial tangent space there, so their equations determine all Taylor derivatives recursively from that value. A zero centre value forces every Taylor coefficient to vanish. The finite mode sums used here are analytic. A scalar fibre has only spin zero; a covector fibre has spins $0,\pm1$; a symmetric rank-two fibre has spins from $-2$ to $2$. A primary with spin outside the corresponding range has zero centre value and therefore zero form factor. Exchange antisymmetry restricts $J$ to odd spin, while exchange symmetry restricts $T$ to even spin. Thus $F$ has only spin zero, $J$ only spin $\pm1$, and the two-identical-scalar stress tensor only spin $0,\pm2$.

Equivalently, this follows by solving the lowering equations with the regular Cartesian centre data. The explicit solutions are \eqref{eq:F}, \eqref{eq:W}, and the conserved scalar-derived tensor used below. A descendant with nonzero orbital angular momentum is not a counterexample: it is not annihilated by both lowering generators.

\section{The annihilation channel}
\label{sec:annihilation}

\subsection{Spin zero}

On free scalar legs the trace identity is
\begin{equation}
 \tau^\mu{}_{\mu}=-\left(\frac14\Box+\mu\right)F_n^{\mathrm{id}}
 =-(C+\mu)F_n^{\mathrm{id}}.
\end{equation}
Define the conserved scalar-to-tensor operator
\begin{equation}
 (\mathscr D S)_{\mu\nu}=
 \nabla_\mu\nabla_\nu S-g_{\mu\nu}(\Box-2)S.
\end{equation}
It obeys $\nabla^\mu(\mathscr D S)_{\mu\nu}=0$ and
$\tr\mathscr D S=-2(\Box-3)S$. Consequently, in the spin-zero primary representation,
\begin{equation}
 \tau_{n0}=\mathscr D S_n,\qquad
 S_n=\frac{C+\mu}{2(4C-3)}F_n^{\mathrm{id}}.
 \label{eq:scalarstress}
\end{equation}
The identity $\mathcal E^{(1)}[g\psi]=-\mathscr D\psi/2$ gives an especially simple response,
\begin{equation}
 q[\tau_{n0}]=-gS_n.
\end{equation}
This computational representative need not be de Donder. Since
\begin{equation}
 \int_\Sigma r\,\dd r\,\dd\varphi\,|F_n^{\mathrm{id}}|^2
 =\frac1{2\pi(2h-1)},
\end{equation}
substitution into \eqref{eq:annpairing} gives
\begin{equation}
 \frac{\gamma^{\mathrm s}_{n0}}G
 =-\frac{2(C+\mu)^2}{(4C-3)(2h-1)}.
\end{equation}
At $n=0,\Delta=3/2$, the ratio in \eqref{eq:scalarstress} is regular after cancellation and the annihilation energy tends to zero. No physical pole is present there.

\subsection{Spin two and higher spin}

The tensor \eqref{eq:taunorm} is transverse and traceless. In that sector
\begin{equation}
 \mathcal E^{(1)}q=-\frac12(\Box+2)q,\qquad
 (\Box+2)\tau_{n2}=4h(h+1)\tau_{n2}.
\end{equation}
Hence $q[\tau_{n2}]=-\tau_{n2}/[4h(h+1)]$. Equations \eqref{eq:annpairing}, \eqref{eq:Wnorm}, and \eqref{eq:taunorm} give
\begin{equation}
 \frac{\gamma^{\mathrm s}_{n2}}G
 =\frac{(n+1)(n+2)(2\Delta+n-1)(2\Delta+n)}
 {(2h-1)(2h+1)(2h+3)}.
\end{equation}
The negative-spin answer is its parity image. Above spin two the primary stress form factor vanishes, so the annihilation-channel coefficient vanishes identically. This is an all-level representation statement, not a truncation of a radial integral.

\section{The crossed-channel Casimir equation}
\label{sec:cross}

\subsection{Collapsing the graviton exchange}

Temporarily distinguish two equal-mass scalar species. Let $x_{kl}$ denote their crossed-exchange coefficient in units of $G$, at chiral levels $(k,l)$; it is symmetric under $k\leftrightarrow l$. The identical-boson crossed contribution is its restriction to $k+l$ even.

On a symmetric rank-two tensor, the sum of the two $SL(2)$ Casimirs is represented by
\begin{equation}
 (\calC q)_{\mu\nu}=\frac12(\Box+6)q_{\mu\nu}-g_{\mu\nu}\tr q.
 \label{eq:tensorcasimir}
\end{equation}
In de Donder gauge the linearized Einstein equation can be written
\begin{equation}
 \mathcal E^{(1)}q=-(\calC-2)\left(q-\frac12g\tr q\right).
\end{equation}
On conserved source pairings this implies
\begin{equation}
 q[T]=-\frac12(\calC-2)^{-1}(T-g\tr T),
 \label{eq:inverse}
\end{equation}
with the prescribed regular reflecting realization and no independently added free metric excitation. The tensor identity \eqref{eq:tensorcasimir} also follows from the homogeneous-bundle Casimir: the symmetric covariant tensor fibre splits into Lorentz spins two and zero, whose isotropy Casimirs are six and zero. Thus $\calC=\Box/2+3P_{\rm traceless}=\Box/2+3-g\tr$. The same relation can be used modulo paired pure-gauge terms, so gauge poles of an auxiliary representative do not create physical poles in the result.

Let $a,b$ be the two incoming scalar legs and $c,d$ the outgoing legs. With
\begin{equation}
 T[a,c]_{\mu\nu}=\nabla_{(\mu}a\nabla_{\nu)}c
 -\frac12g_{\mu\nu}(\nabla a\cdot\nabla c+\mu ac),
\end{equation}
the crossed energy pairing has normalization $-2\kap^2\int T[a,c]q[T[b,d]]$. Acting with the crossed Casimir minus two therefore collapses the propagator to the local vertex
\begin{equation}
 \kap^2\mathcal K(a,b;c,d),\qquad
 \mathcal K=T[a,c]^{\mu\nu}T[b,d]_{\mu\nu}
 -\tr T[a,c]\tr T[b,d].
 \label{eq:contactK}
\end{equation}
This identity follows from \eqref{eq:inverse} and fixes the contact term; there is no adjustable low-spin parameter.

\subsection{The contact source in a primary basis}

For a normalized incoming primary, set $F=\bra0\phi_1\phi_2\ket P$,
$J=\bra0\phi_1\dd\phi_2-\phi_2\dd\phi_1\ket P$, and
$K_F=\frac12(\Box-2\mu)F$. Expanding \eqref{eq:contactK} and pairing incoming and outgoing primary wavefunctions gives
\begin{align}
 \int_\Sigma\mathcal K
 =\int_\Sigma\left\{
 \frac12|K_F|^2-\frac1{16}|\dd J|^2
 -\frac\mu4\bigl(\nabla F\cdot\nabla F^*+J\cdot J^*\bigr)
 -\frac32\mu^2|F|^2\right\}.
 \label{eq:Kdecomp}
\end{align}
Here and below $\int_\Sigma$ includes the measure $r\dd r\dd\varphi$ and the resonant time average. Integration by parts is legitimate for these normalizable primary form factors; the time endpoint terms cancel in the pairing.

For spin zero, \eqref{eq:Kdecomp} reduces to
$(C-\mu)(2C+\mu)\int|F_n|^2$. For spin one, $F=0$ and
$\nabla^\mu(\dd J)_{\mu\nu}=4h^2J_\nu$, so it reduces to
$(2h^2-\mu)\int J_n\cdot J_n^*/4$. The form-factor norms are
\begin{equation}
 \int|F_n|^2=\frac1{4\pi(2h-1)},\qquad
 \int J_n\cdot J_n^*=\frac{h(n+1)(2\Delta+n-1)}{\pi(2h-1)(2h+1)}.
\end{equation}
Thus the collapsed contact coefficients are
\begin{align}
 S_{nn}&=\frac{4(C-\mu)(2C+\mu)}{2h-1},\label{eq:S0}\\
 S_{n+1,n}=S_{n,n+1}
 &=\frac{4h(n+1)(2\Delta+n-1)(2h^2-\mu)}{(2h-1)(2h+1)},\label{eq:S1}\\
 S_{kl}&=0,\qquad |k-l|\geq2.\label{eq:Shigh}
\end{align}
The finite support of this contact source is particularly useful. It is not the same spin support as the annihilation contribution.

\subsection{A three-term Casimir action}

Use the chiral blocks and free coefficients
\begin{equation}
 k_h(z)=z^h\,{}_2F_1(h,h;2h;z),\qquad
 p_k=\frac{(\Delta)_k^2}{k!(2\Delta+k-1)_k}.
\end{equation}
The free distinguishable two-particle kernel is
\begin{equation}
 \left(\frac{z\bar z}{(1-z)(1-\bar z)}\right)^\Delta
 =\sum_{k,l\geq0}p_kp_l\,k_{\Delta+k}(z)k_{\Delta+l}(\bar z).
\end{equation}
At first order the coefficient of $\log(z\bar z)$ is
$\frac G2\sum p_kp_l x_{kl}k_{\Delta+k}k_{\Delta+l}$. The crossed chiral Casimir is
\begin{equation}
 D_t=z(1-z)^2\partial_z^2+(z-1)(2\Delta+z-1)\partial_z
 +\frac{\Delta^2}{z}-\Delta.
\end{equation}
It acts tridiagonally:
\begin{align}
 D_tk_h={}&(\Delta-h)^2k_{h-1}
 +\frac{\mu-h(h-1)}2k_h
 +\frac{h^2(\Delta+h-1)^2}{4(2h-1)(2h+1)}k_{h+1}.
 \label{eq:threeblock}
\end{align}
Appendix~\ref{app:certificate} proves this identity by comparing the coefficient with an arbitrary hypergeometric series index, not just the first few coefficients.

After including the weights $p_k$, define
\begin{align}
 a_k&=\frac{(k+1)(\Delta+k)(2\Delta+k-1)}{2(2\Delta+2k-1)},&
 c_k&=\frac{k(\Delta+k-1)(2\Delta+k-2)}{2(2\Delta+2k-1)},\\
 b_k&=\frac{\mu-(\Delta+k)(\Delta+k-1)}2=-a_k-c_k,&
 (\mathcal L f)_k&=a_kf_{k+1}+b_kf_k+c_kf_{k-1}.
 \label{eq:abc}
\end{align}
Since $c_0=0$, no negative-level datum is required. Taking the logarithmic coefficient of the collapsed exchange identity gives the exact equation
\begin{equation}
 \boxed{(\mathcal L_k+\mathcal L_l-2)x_{kl}=S_{kl}.}
 \label{eq:crossrec}
\end{equation}
Differentiating the logarithm produces nonlogarithmic terms, which do not enter this equation. This is the usual AdS representation/block relation; it can be used as a harmonic-analysis device without assuming a separately constructed nonperturbative boundary CFT.

\subsection{A solution and the required high-spin input}

The known crossed stress-tensor exchange gives, for $|k-l|>2$,
\begin{equation}
 x_{kl}=-\frac{12}{cG}\left[
 C_2\bigl(\Delta+\min(k,l)\bigr)-2C_2(\Delta/2)\right]
 =U_{\min(k,l)},\qquad c=\frac3{2G}+O(1).
 \label{eq:tail}
\end{equation}
We use Ref.~\cite{KSS}, Section 4.2, only in this above-exchanged-spin range. In particular, \eqref{eq:tail} is not an assumed formula for spin zero, one, or two.

The following extension solves \eqref{eq:crossrec} at every level:
\begin{equation}
 x_{kl}=
 \begin{cases}
 U_n\dfrac{2\Delta+2n-2}{2\Delta+2n-1},&k=l=n,\\[6pt]
 U_{\min(k,l)},&k\ne l.
 \end{cases}
 \label{eq:xsolution}
\end{equation}
For $|k-l|\geq2$, verification reduces to $\mathcal L1=0$ and $\mathcal LU=2U$. On the two remaining bands, direct rational substitution gives precisely \eqref{eq:S0} and \eqref{eq:S1}. This proves existence for arbitrary $\Delta,n$; uniqueness is a separate issue addressed next.

\subsection{All-level uniqueness of the low-spin completion}
\label{sec:uniqueness}

Suppose $e_{kl}=e_{lk}$ is the difference of two solutions with the same contact source and the same tail \eqref{eq:tail}. Then $e_{kl}=0$ for $|k-l|\geq3$ and
$(\mathcal L_k+\mathcal L_l-2)e_{kl}=0$. Denote its three possibly nonzero bands by
\begin{equation}
 X_n=e_{n+2,n},\qquad Y_n=e_{n+1,n},\qquad Z_n=e_{n,n}.
\end{equation}
The equations on gaps three, two, and zero give
\begin{align}
 X_{n+1}&=-\frac{c_{n+3}}{a_n}X_n,\label{eq:uniqueX}\\
 Y_{n+1}&=-\frac{c_{n+2}Y_n+(b_{n+2}+b_n-2)X_n}{a_n},\label{eq:uniqueY}\\
 Z_n&=\frac{a_nY_n+c_nY_{n-1}}{1-b_n}.\label{eq:uniqueZ}
\end{align}
At $n=0$ the term with $c_0$ is omitted. All denominators are nonzero for $\Delta>1$. Every coefficient is therefore determined by just two numbers, $(Y_0,X_0)$.

The remaining gap-one equations are
\begin{equation}
 R_n=c_{n+1}Z_n+a_nZ_{n+1}+(b_{n+1}+b_n-2)Y_n
 +a_{n+1}X_n+c_nX_{n-1}=0.
 \label{eq:uniqueR}
\end{equation}
Substituting \eqref{eq:uniqueX}--\eqref{eq:uniqueZ} into $R_0,R_1$ yields a $2\times2$ homogeneous system $M(\Delta)(Y_0,X_0)^T=0$. Its determinant is
\begin{equation}
 \boxed{\det M(\Delta)=
 \frac{960(\Delta+1)^2
 (3\Delta^5+24\Delta^4+63\Delta^3+68\Delta^2+33\Delta+6)}
 {\Delta(2\Delta+3)(2\Delta+5)(3\Delta+2)(5\Delta+4)}>0.}
 \label{eq:det}
\end{equation}
Thus $X_0=Y_0=0$ and all three bands vanish. This proves uniqueness for every $n$; no growth assumption, finite-level fit, or infinite-mode completion is needed for this finite-band argument. Restriction of \eqref{eq:xsolution} to the even-spin identical-boson sector proves \eqref{eq:crossanswer}. Adding Section~\ref{sec:annihilation} proves Theorem~\ref{thm:main}.

\section{Independent canonical check from circular data}
\label{sec:circular}

The preceding proof bypasses the difficult direct summation of general radial Hahn-polynomial compressions. Those compressions provide an independent check because they start from the exact circular Einstein constraints, rather than the crossed-channel equation.

In polar areal gauge, put
\begin{equation}
 \dd s^2=-Fe^{-2D}\dd t^2+\frac{\dd r^2}{F}+r^2\dd\varphi^2,
 \qquad F=f-\kap^2M,\qquad \Pi=\frac{e^D}{F}\dot\phi.
\end{equation}
The reduced form, mass constraint, and lapse equation are
\begin{align}
 \Omega_{\mathrm{red}}&=2\pi\int_0^\infty r\,\delta\Pi\wedge\delta\phi\,\dd r,\\
 M'&=\frac r2\left[F(\Pi^2+\phi'^2)+\mu\phi^2\right],&
 D'&=-\frac{\kap^2r}{2}(\Pi^2+\phi'^2).
\end{align}
Regularity and boundary time fix $M(0)=0$ and $D(\infty)=0$. The vacuum-subtracted Hamiltonian is $H=2\pi M(\infty)$. With $X=\Pi^2+\phi'^2$,
\begin{align}
 H&=\pi\int_0^\infty\!r(fX+\mu\phi^2)
 \exp\!\left[-\frac{\kap^2}{2}\int_r^\infty\!sX(s)\dd s\right]\dd r,\\
 H&=H_0+\kap^2H_4+O(\kap^4),&
 H_4&=-\pi\int_0^\infty rX(r)M_0(r)\dd r,\label{eq:H4}\\
 M_0(r)&=\frac12\int_0^r s\,[f(s)X(s)+\mu\phi(s)^2]\dd s.
\end{align}
The gravitational momentum term vanishes in this circular reduction: $K_{\varphi\varphi}=0$ implies $p^{rr}=0$, while $\delta\gamma_{\varphi\varphi}=0$. Thus the stated scalar Cauchy coordinates are canonical without an extra perturbative Darboux correction.

For the lowest mode $u=f^{-\Delta/2}/\sqrt{2\pi}$, the coefficient of $b^2\bar b^2$ in $H_4$ is
\begin{equation}
 C_{00}=\frac{\Delta^2(7+2\Delta-8\Delta^2)}{16\pi(4\Delta^2-1)}.
\end{equation}
The normalized two-boson expectation is $2\kap^2C_{00}$, agreeing with \eqref{eq:full0}. The nonresonant quartic monomials do not shift this isolated level at first order. Adding a separate exchange term to \eqref{eq:H4} would double count the constraints already integrated out.

At general circular level $N$, use normalized unordered pairs
$b^\dagger_{r,0}b^\dagger_{N-r,0}\ket0/\sqrt{1+\delta_{2r,N}}$ and
\begin{equation}
 P_i(x)=P_i^{(\Delta-1,0)}(1-2x),\qquad x=f^{-1}.
\end{equation}
The six four-leg pairings in \eqref{eq:H4} reduce to finite Mellin moments of Jacobi polynomials. Same-sign pairings isolate the annihilation channel, and mixed-sign pairings isolate crossed exchange. Their compression is not a closed angular-momentum Hamiltonian block. To extract primary coefficients, use the orthogonal chiral coupling matrix $U^{(N)}$ obtained from Hahn polynomials:
\begin{equation}
 R^{(N)}_{rs}=\sum_{\substack{0\leq k,l\leq N\\k+l\text{ even}}}
 s_rs_sU_{rk}^{(N)}U_{rl}^{(N)}U_{sk}^{(N)}U_{sl}^{(N)}
 \gamma_{\min(k,l),|k-l|},\quad
 s_r=\sqrt{\frac2{1+\delta_{2r,N}}}.
 \label{eq:compression}
\end{equation}
The top Hahn column never vanishes for $\Delta>1$. Subtracting all earlier levels leaves an invertible finite reconstruction problem. The supplied implementation tests the unused off-diagonal residuals rather than discarding them. The detailed finite-sum implementation and its provenance are included with this draft \cite{notes}.

For $N=n+|\ell|\leq12$, the independent exact-rational calculation evaluates 49 primary coefficients for each of
$\Delta=21/20,3/2,2,7/3$. All three channel choices (total, crossed, annihilation) agree with \eqref{eq:crossanswer}--\eqref{eq:sanswer}, giving 588 exact coefficient comparisons. Orthogonality and unused off-diagonal residuals vanish in every run. This is finite computational evidence distinct from the all-level analytic argument in Sections~\ref{sec:annihilation} and \ref{sec:cross}.

\begin{table}[ht]
\centering
\begin{tabular}{rrrrr}
\toprule
$n$ & $|\ell|$ & $E_0$ at $\Delta=2$ & $g^{\mathrm x}_{n\ell}$ & $g_{n\ell}$\\
\midrule
0&0&4&$-32/3$&$-56/5$\\
1&0&6&$-192/5$&$-1368/35$\\
0&2&6&$-16$&$-552/35$\\
2&0&8&$-576/7$&$-416/5$\\
1&2&8&$-48$&$-1000/21$\\
0&4&8&$-16$&$-16$\\
3&0&10&$-1280/9$&$-11040/77$\\
2&2&10&$-96$&$-7352/77$\\
1&4&10&$-48$&$-48$\\
0&6&10&$-16$&$-16$\\
\bottomrule
\end{tabular}
\caption{Selected coefficients. Each positive spin has an equal-energy parity partner. Circular product-state eigenvalues alone would not produce this primary table.}
\label{tab:spectrum}
\end{table}

\section{Interpretation and remaining physical scope}

\subsection{Binding, low-spin effects, and special masses}

The ground coefficient simplifies to
\begin{equation}
 g_{00}=\frac{2\Delta^2(7+2\Delta-8\Delta^2)}{4\Delta^2-1}.
\end{equation}
It is positive for $1<\Delta<(1+\sqrt{57})/8$ and negative above that interval. This does not follow from a universal attractive-potential argument: when $1<\Delta<2$, $\mu<0$ and the finite-radius mass $M_0(r)$ in \eqref{eq:H4} is not pointwise positive. Positivity of the total energy does not make every partial radial mass positive. The correct statement is a positive connected shift of this particular two-particle state in the specified theory, not a general claim that gravity becomes repulsive.

At fixed $n$ and large $\Delta$ the leading shift is attractive. At large radial level with fixed $\Delta$, the universal term grows quadratically while the low-spin corrections grow linearly. These asymptotics describe the coefficients; perturbation theory still requires the retained states to remain in the weak-backreaction regime.

The auxiliary no-log de Donder construction in the source notes excludes
$\Delta_*=(1+\sqrt5)/2$. No denominator in the physical coefficient formulas vanishes there. The conformal response used in the spin-zero calculation and the canonical reduction show why a gauge-indicial restriction should not be identified with a singular energy. The formulas extend there by continuity at each fixed level. The only apparent denominator zero inside $\Delta>1$ is the removable $n=0,\Delta=3/2$ point already discussed.

\subsection{Boundary gravitons and degeneracies}

Physical Brown--Henneaux excitations remain in the theory. A non-null descendant of a primary of gap $E_P$ has gap $E_P+N_L+N_R$; boundary-graviton states contribute to the multiplicities of these modules. The present calculation fixes the connected scalar-primary binding coefficient, not the complete change of basis from every bare scalar/graviton Fock state to an interacting module basis.

In particular, the scalar-only circular compression must not be presented as the entire physical Hilbert space. A separate analysis of the full projected Hamiltonian is required wherever a proposed spectral claim depends on mixing with additional physical sectors. The exchange calculation already includes the virtual metric response, so that analysis must not add a second copy of the same exchange. A useful further check of the underlying canonical construction is a directly evaluated noncircular matrix element; the all-level coefficient proof above does not by itself supply that full Fock-space dressing map.

\subsection{What is and is not established}

The theorem gives a closed all-$n$ completion of the minimal connected tree result, using a known high-spin tail, explicitly normalized primary form factors, and an all-level uniqueness argument. It does not follow merely from fitting the previous finite radial data. Conversely, finite exact checks of the radial recursion do not establish a theorem about infinite sums of modes. No UV regulator-specific self-energy, independent contact coupling, higher-loop term, nonstandard scalar quantization, or infinite-mode convergence statement has been included.

\appendix
\section{Normalization of the local form factors}
\label{app:normalization}

The norm identity underlying \eqref{eq:CG} is
\begin{equation}
 \sum_{p=0}^K\frac{\binom Kp}{(\Delta)_p(\Delta)_{K-p}}
 =\frac{(2\Delta+K-1)_K}{(\Delta)_K^2}.
\end{equation}
At the centre only scalar modes of angular momentum zero contribute to $F$, and only angular momentum one contributes to a first spatial derivative. With $w=x+iy$ a regular complex Cartesian coordinate, \eqref{eq:centremode} gives
\begin{equation}
 \dd u_{r+1,r}\big|_0=\frac{(-1)^r}{\sqrt{2\pi}}
 \sqrt{(r+1)(\Delta+r)}\,\dd w.
\end{equation}
Set
\begin{align}
 R_1^2&=\frac{h(n+1)(2\Delta+n-1)}{2(2h-1)},\\
 R_2^2&=\frac{h(h+1)(n+1)(n+2)(2\Delta+n-1)(2\Delta+n)}{4(2h-1)(2h+1)}.
\end{align}
The coefficient formula \eqref{eq:CG} immediately gives the $r$-independent ratios
\begin{align}
 \frac{v_r^{(n+1)}}{v_r^{(n)}}
 \sqrt{(n-r+1)(\Delta+n-r)}&=R_1,\\
 \frac{v_{r+1}^{(n+1)}}{v_r^{(n)}}
 \sqrt{(r+1)(\Delta+r)}&=-R_1,\\
 \frac{v_{r+1}^{(n+2)}}{v_r^{(n)}}
 \sqrt{(r+1)(\Delta+r)(n-r+1)(\Delta+n-r)}&=-R_2.
\end{align}
Consequently the current centre sum is $(-1)^nR_1\dd w/\pi$, while the identical stress centre sum is
$(-1)^{n+1}R_2\dd w^2/(\sqrt2\pi)$.
Since $W^{(E,1)}|_0=i\dd w$ and
$W^{(E,2)}|_0=-\dd w^2$, this proves
$|A_1|^2=R_1^2/\pi^2$ and $|A_2|^2=R_2^2/(2\pi^2)$.

The elementary contraction
\begin{equation}
 v\cdot v=0,\qquad v\cdot v^*=\frac2{r^2f}
\end{equation}
gives \eqref{eq:Wnorm} by integrating
$2\pi 2^s\int_0^\infty r f^{-E-s}\dd r$.
For completeness, the two lowering vectors used to check the primary equations are
\begin{equation}
 D_\sigma=\frac12 e^{it-i\sigma\varphi}
 \left(\frac{ir}{\sqrt f}\partial_t+\sqrt f\partial_r
 -\frac{i\sigma\sqrt f}{r}\partial_\varphi\right),\qquad \sigma=\pm1.
\end{equation}
The supplied tensor script checks $\mathcal L_{D_\sigma}W=0$, transversality, tracelessness where appropriate, and every component of \eqref{eq:Weigen}, with symbolic $E$ and $r$.

\section{Block identities and arbitrary-index certificates}
\label{app:certificate}

A useful check on the all-level high-spin coefficient extraction is the identity
\begin{equation}
 \sum_{n\geq0}p_n U_n k_{\Delta+n}(z)
 =-4\Delta^2\frac{1+z}{1-z}\left(\frac{z}{1-z}\right)^\Delta.
 \label{eq:tailidentity}
\end{equation}
Indeed, the direct-channel chiral Casimir
$D_s=z^2(1-z)\partial_z^2-z^2\partial_z$ has eigenvalue $h(h-1)$ on $k_h$.
Acting with $4\mu-8D_s$ on the free expansion of
$F(z)=[z/(1-z)]^\Delta$ proves \eqref{eq:tailidentity}, since
$D_sF=[\Delta^2/(1-z)-\Delta]F$.
The rational factor on the right is precisely the logarithmic factor in crossed stress-tensor exchange. The use of inversion above spin two remains the input specified in Section~\ref{sec:cross}; this identity does not extend inversion to spin zero by itself.

Write $k_h(z)=z^h\sum_{j\geq0}f_j(h)z^j$, where
$f_j(h)=(h)_j^2/[(2h)_j j!]$. For $j\geq2$, define
\begin{align}
 r_1&=\frac{f_{j-1}(h)}{f_j(h)}
 =\frac{j(2h+j-1)}{(h+j-1)^2},\\
 r_2&=\frac{f_{j-2}(h)}{f_j(h)}
 =\frac{j(j-1)(2h+j-1)(2h+j-2)}{(h+j-1)^2(h+j-2)^2},\\
 r_-&=\frac{f_j(h-1)}{f_j(h)}
 =\frac{(h-1)^2(2h+j-2)(2h+j-1)}{(h+j-1)^2(2h-2)(2h-1)},\\
 r_+&=\frac{f_{j-2}(h+1)}{f_j(h)}
 =\frac{j(j-1)(2h)(2h+1)}{h^2(h+j-1)^2}.
\end{align}
The coefficient of $z^{h-1+j}$ in the left side of \eqref{eq:threeblock}, divided by $f_j(h)$, is
\begin{equation}
 (h+j-\Delta)^2+r_1[-2(h+j-1)(h+j-1-\Delta)-\Delta]
 +r_2(h+j-2)^2.
\end{equation}
The right side, divided by the same coefficient, is
\begin{equation}
 (\Delta-h)^2r_-+\frac{\mu-h(h-1)}2r_1
 +\frac{h^2(\Delta+h-1)^2}{4(2h-1)(2h+1)}r_+.
\end{equation}
Clearing denominators gives an identically zero polynomial in $j,h,\Delta$. The $j=0,1$ coefficients agree directly. This proves the entire three-term identity. The symbolic script includes this arbitrary-$j$ certificate and the determinant calculation \eqref{eq:det}; its checks do not replace the field-theoretic derivation of the contact source.

\section{Reproducibility and provenance}

The starting canonical reduction and radial reconstruction are recorded in the source notes at the pinned repository commit \cite{notes}. The present draft adds the channel decomposition, local form-factor evaluation, crossed contact source, finite-band uniqueness proof, and analytic closed formulas. The calculations supplied with this draft require Python 3 and SymPy; the radial script itself uses only the standard library.

From the source directory, run
\begin{verbatim}
python verify_symbolic.py
python verify_tensor_modes.py
python verify_radial.py --level 12 --output radial_report.json
pdflatex -interaction=nonstopmode -halt-on-error paper.tex
pdflatex -interaction=nonstopmode -halt-on-error paper.tex
\end{verbatim}
The first script checks 21 rational identities with symbolic parameters, including the arbitrary hypergeometric index and the low-spin uniqueness determinant. The second checks the spin-one and spin-two tensor fields component by component. The third independently reconstructs 49 coefficients at each of four rational masses in three channels. Its exact arithmetic output is retained in \texttt{radial\_report.json}. The arbitrary-level statements in the text rest on the analytic arguments, not on increasing the numerical cutoff.

This is a working research draft. The high-spin input is explicitly cited, and no claim of an exhaustive novelty search is made. The full scalar/boundary-graviton mixing map, higher loops, and infinite-mode functional analysis remain outside its scope.

\begin{thebibliography}{9}
\bibitem{FS}
A.~L.~Fitzpatrick and D.~Shih,
\emph{Anomalous Dimensions of Non-Chiral Operators from AdS/CFT},
\arxiv{1104.5013}.

\bibitem{KSS}
P.~Kraus, A.~Sivaramakrishnan and R.~Snively,
\emph{Late time Wilson lines},
\arxiv{1810.01439}, especially Section 4.2.

\bibitem{CH}
S.~Caron-Huot,
\emph{Analyticity in Spin in Conformal Theories},
\arxiv{1703.00278}.

\bibitem{ABP}
L.~F.~Alday, A.~Bissi and E.~Perlmutter,
\emph{Holographic Reconstruction of AdS Exchanges from Crossing Symmetry},
\arxiv{1705.02318}.

\bibitem{HW}
D.~Harlow and J.-q.~Wu,
\emph{Covariant phase space with boundaries},
\arxiv{1906.08616}.

\bibitem{notes}
GaoZ1en, \emph{Einstein--scalar response and particle-spectrum research notes},
\texttt{obsidian\_note}, commit
\href{https://github.com/GaoZ1en/obsidian_note/commit/3ea676ef5a43cd11a18f551b0adcac222abc8b9a}{\texttt{3ea676e}} (12 September 2026),
\texttt{Articles/Quantization in AdS/linearized gravity/}.
In particular, \emph{gravity scalar one and two particle spectrum.md} and
\texttt{scripts/gravity\_scalar\_radial\_blocks.wl}.
\end{thebibliography}
\end{document}
